\documentclass[11pt]{article}
\usepackage[margin=1.1in]{geometry}
\usepackage{amsmath,amssymb,amsthm}
\usepackage{enumitem}
\usepackage[colorlinks=true,linkcolor=blue,citecolor=blue,urlcolor=blue]{hyperref}

\newtheorem{theorem}{Theorem}[section]
\newtheorem{lemma}[theorem]{Lemma}
\newtheorem{proposition}[theorem]{Proposition}
\newtheorem{corollary}[theorem]{Corollary}
\newtheorem{remark}[theorem]{Remark}
\theoremstyle{definition}
\newtheorem{definition}[theorem]{Definition}
\newtheorem{condition}[theorem]{Condition}

\newcommand{\pP}{\textbf{[P]}}
\newcommand{\pC}{\textbf{[C]}}
\newcommand{\pO}{\textbf{[O]}}
\newcommand{\Tr}{\operatorname{Tr}}

\title{\textbf{Spin-Current Sources and Ball Modular Hamiltonians:\\
Resolving the Torsionful First Law}}
\author{Robert W.\ McGwier\\
\small CTO and Staff Scientist, Cohere Technology Group, Herndon, VA}
\date{August 7, 2026 --- Version 2}

\begin{document}
\maketitle

\begin{abstract}
\noindent
In the companion synthesis \cite{McGwier2026} it was proved that, at linear
order about the vacuum, the entanglement first law is equivalent to the
linearized Einstein--Cartan equations in the metric sector, with torsion
fixed algebraically by the spin density of the state; the open remainder was
the status of the first law when the boundary theory carries an
\emph{explicit source} for the spin current. Here we resolve that remainder
into a theorem sector and a counterexample sector. For minimally coupled
Dirac matter in $d=4$: (i) only the totally antisymmetric part of a
spin-current source couples, and it is equivalent to a background axial
vector field $\beta_\mu$ coupled to $J_5^\mu$ (Reduction Lemma, derived);
(ii) at first order on a flat, otherwise sourceless background, ball modular
Hamiltonians depend on $\beta$ only through its field strength
$d\beta$---pure-gauge torsion sources are unitarily trivial modulo the
anomaly (Gauge Lemma, derived); (iii) the deformed modular Hamiltonian is
governed by the Faulkner--Leigh--Parrikar--Wang / S\'arosi--Ugajin
first-order formula, whose nonlocal remainder vanishes on null planes by the
Markov property of the vacuum, yielding an unconditional first-law/EC
equivalence in the null-plane limit and a conditional one for finite balls,
with the required condition stated exactly. (iv) Mixed-symmetry torsion
sources coupling through non-minimal matter are relevant deformations of
dimension $d-1$: they destroy the conformal frame on which the ball
construction rests, and the naive torsionful first law is \emph{false} in
that sector. Every claim is labeled \pP{}/\pC{}/\pO{}.
\end{abstract}

\section{The Problem, Stated Exactly}

Proposition 6.3.2 of \cite{McGwier2026} established: about a maximally
symmetric vacuum, the first law $\delta S_B=\delta\langle H_B\rangle$ on all
balls is equivalent to the linearized Einstein--Cartan (EC) equations, with
the torsion sector fixed algebraically by the spin density of the perturbed
state. Its Remark 6.3.3(i) recorded the open remainder: the case of an
\emph{explicit source} $b_{\lambda\mu\nu}(x)$ coupled to the spin current,
\begin{equation}
S \;\longrightarrow\; S \;+\; \int d^dx\;
b_{\lambda\mu\nu}(x)\, s^{\lambda\mu\nu}(x)\,,
\label{eq:deformation}
\end{equation}
which holographically corresponds to a torsionful boundary condition
\cite{GallegosGursoy2020,GGY2021}. The question: does the first law, in the
theory deformed by \eqref{eq:deformation}, remain equivalent to linearized
EC gravity with contortion boundary data set by $b$? We work at first order
in $b$ and first order in the state perturbation $\lambda$, tracking the
joint order $O(b\lambda)$ where the two expansions couple. Boundary
dimension $d=4$ (bulk AdS$_5$) unless stated.

\section{Reduction Lemma: Torsion Sources are Axial Sources}

\begin{lemma}[Reduction]
\label{lem:reduction}
For a minimally coupled massless Dirac field in $d=4$, the canonical spin
current is totally antisymmetric and dual to the axial current: with the
conventions of \cite{HehlDatta1971,HehlRMP1976},
\begin{equation}
s^{\lambda\mu\nu} \;=\; -\tfrac{1}{4}\,\epsilon^{\lambda\mu\nu\rho}\,
J_{5\rho}\,,\qquad J_5^\rho=\bar\psi\gamma^5\gamma^\rho\psi\,.
\label{eq:dualspin}
\end{equation}
Consequently the deformation \eqref{eq:deformation} reduces to
\begin{equation}
\int d^4x\; b_{\lambda\mu\nu}\,s^{\lambda\mu\nu}
\;=\;\int d^4x\;\beta_\rho(x)\,J_5^\rho(x)\,,\qquad
\beta^\rho \;\equiv\; -\tfrac{1}{4}\,\epsilon^{\rho\lambda\mu\nu}\,
b_{[\lambda\mu\nu]}\,,
\label{eq:axialreduction}
\end{equation}
and the vector and mixed-symmetry irreducible parts of
$b_{\lambda\mu\nu}$ decouple identically.
\end{lemma}

\begin{proof}
The canonical spin current of the Dirac Lagrangian is
$s^{\lambda\mu\nu}=\tfrac{1}{8}\bar\psi\,
\{\gamma^\lambda,[\gamma^\mu,\gamma^\nu]\}\,\psi$. The Clifford identity
$\{\gamma^\lambda,[\gamma^\mu,\gamma^\nu]\}\propto
\epsilon^{\lambda\mu\nu\rho}\gamma^5\gamma_\rho$ (the anticommutator
projects onto the totally antisymmetrized product
$\gamma^{[\lambda}\gamma^\mu\gamma^{\nu]}$, which in $d=4$ is dual to
$\gamma^5\gamma_\rho$) gives \eqref{eq:dualspin}. Total antisymmetry means
the contraction in \eqref{eq:deformation} projects $b$ onto its totally
antisymmetric part; the remaining irreps of a rank-3 tensor with the index
symmetry of torsion (the trace-vector and the mixed-symmetry tensor part)
contract to zero. \pP{}
\end{proof}

\begin{remark}
This is the boundary shadow of the bulk Hehl--Datta structure
\cite{HehlDatta1971}: for minimal Dirac matter only the axial (totally
antisymmetric) piece of contortion is dynamical, and the induced contact
interaction is axial--axial. The Lemma states that the same projection
governs which torsion \emph{sources} act on the boundary theory. Duality
\eqref{eq:dualspin} is specific to $d=4$; total antisymmetry of the Dirac
spin current holds in all $d$. \pP{}
\end{remark}

\section{Gauge Lemma: Only the Curl of the Source Matters}

\begin{lemma}[Gauge]
\label{lem:gauge}
Let the boundary CFT be massless Dirac matter on flat $\mathbb{R}^{1,3}$
with no other background fields, deformed by
$\int\beta\cdot J_5$. If $\beta_\mu=\partial_\mu\alpha$, then at first
order in $\beta$ the deformation is implemented by the unitary chiral
rotation $U_\alpha=\exp\!\big(i\!\int
\alpha\,j^0_5\big)$-type conjugation, and every ball modular Hamiltonian
transforms by conjugation:
$H_B^{(\beta)}=U_\alpha\,H_B\,U_\alpha^\dagger+O(\beta^2)$. Hence all
first-order modular data---in particular relative entropies and the first
law---depend on $\beta$ only through $d\beta$.
\end{lemma}

\begin{proof}
$\int\partial_\mu\alpha\,J_5^\mu=-\int\alpha\,\partial_\mu J_5^\mu$ up to a
surface term (take $\alpha$ compactly supported). For massless Dirac,
$\partial_\mu J_5^\mu$ vanishes as an operator equation up to the anomaly,
$\partial\!\cdot\!J_5\propto F\tilde F$ plus curvature terms; on flat space
with no gauge background these insertions are themselves $O(\beta)$ or
higher, so the deformation vanishes at strict first order and the chiral
rotation implementing it has trivial Jacobian at this order (Fujikawa
measure contributions enter only multiplied by background field strengths).
Conjugation of modular operators under unitaries is the general identity
$\rho\to U\rho U^\dagger \Rightarrow K\to UKU^\dagger$. \pP{}, with the
anomaly caveat stated: at $O(\beta^2)$ and in the presence of gauge or
curvature backgrounds, the Nieh--Yan/torsional anomaly channel
\cite{NiehYan1982,ChandiaZanelli1997} activates and pure-gauge triviality
fails. \pO{} beyond first order.
\end{proof}

\begin{corollary}
Constant torsion sources (homogeneous $\beta_\mu$, e.g.\ a chiral chemical
potential) are \emph{not} pure gauge on $\mathbb{R}^{1,3}$ with compact
support removed---$\beta=$ const has $d\beta=0$ but is not globally exact
with the required falloff---so the Lemma constrains local physics only:
local first-order modular response is a functional of $d\beta$ and of the
holonomy data of $\beta$. \pP{}
\end{corollary}

\section{The Deformed Modular Hamiltonian}

Let $K_B$ be the vacuum modular Hamiltonian of ball $B$ in the undeformed
CFT (the CHM form \cite{CHM2011}), and let
$O_\beta=\int\beta\cdot J_5$. The first-order shift of the modular
Hamiltonian under a perturbation of the state or of the theory is given by
the universal modular-smearing formula of perturbative modular theory
\cite{FLPW2016,SarosiUgajin2018}: $\delta K_B$ is the vacuum-modular-flow
smearing of the induced density-matrix perturbation with kernel
$\big(4\sinh^2\!\tfrac{s+i\epsilon}{2}\big)^{-1}$. We do not require the
kernel's normalization; we require only three structural facts, each
established in the cited literature:

\begin{enumerate}[itemsep=2pt]
\item \textbf{(Linearity)} $\delta K_B$ is linear in $\beta$. \pP{}
\item \textbf{(Local term)} For the Rindler wedge $W=\{x^1>|t|\}$, the
deformation contributes the boost-weighted local term
\begin{equation}
\delta K_W^{\mathrm{loc}}
=2\pi\!\int_{t=0,\,x^1>0}\!\! d^{3}x\;x^1\,
\beta_\mu(x)\,J_5^\mu(x)\,,
\label{eq:localterm}
\end{equation}
the direct analogue of the FLPW result that first-order deformations enter
the wedge modular Hamiltonian with boost weight \cite{FLPW2016}. \pP{}
\item \textbf{(Null-plane collapse)} In the limit in which the entangling
surface is deformed onto a null plane, the nonlocal remainder
$\delta K^{\mathrm{nl}}\equiv\delta K-\delta K^{\mathrm{loc}}$ vanishes:
the vacuum restricted to null-plane algebras is a Markov state, modular
Hamiltonians of null cuts are additive and exactly local
\cite{CasiniTesteTorroba2017}, leaving no room for a nonlocal piece.
\pP{}
\end{enumerate}

For finite balls, $\delta K_B^{\mathrm{nl}}\neq 0$ in general; its exact
form is the modular smearing of $O_\beta$ minus the local term. This
motivates:

\begin{condition}[NL]
\label{cond:NL}
For the class of state perturbations $|\delta\Psi\rangle$ dual to
linearized bulk EC solutions, the cross expectation
$\langle 0|\delta K_B^{\mathrm{nl}}|\delta\Psi\rangle+\mathrm{c.c.}$
vanishes for all balls $B$ at joint order $O(\beta\lambda)$.
\end{condition}

\section{Main Results}

\begin{theorem}[Null-plane torsionful first law]
\label{thm:nullplane}
For massless minimally coupled Dirac matter in $d=4$ deformed by a smooth
compactly supported spin-current source, the entanglement first law on null
cuts is equivalent, at joint order $O(\beta\lambda)$, to the linearized
Einstein--Cartan equations integrated against the corresponding null
generators, with contortion boundary data given by the axial reduction of
$b$ (Lemma \ref{lem:reduction}).
\end{theorem}

\begin{proof}[Derivation]
On null planes the modular Hamiltonian of the deformed theory is exactly
local with boost weight (facts 2--3 above; \cite{FLPW2016,
CasiniTesteTorroba2017}), so
$\delta\langle K^{(\beta)}\rangle$ is the null integral of
$\delta\langle T_{++}\rangle+\beta\!\cdot\!\delta\langle J_5\rangle$-type
insertions with known weights. On the bulk side, the null-plane limit of
the FGHMV form argument \cite{Faulkner2014} reduces the RT variation to the
integrated null-null linearized equation of motion; by the decomposition
proposition of \cite{McGwier2026}~\S6.3, the EC system at this order is the
Einstein sector sourced by the Belinfante tensor of the $\beta$-deformed
theory plus the algebraic Cartan sector, whose content is exactly the axial
identification of Lemma \ref{lem:reduction}. Matching the two sides is then
the FGHMV null matching with the shifted source. The anomaly channel
\cite{NiehYan1982,ChandiaZanelli1997} is $O(\beta^2)$ and does not enter.
\pP{} at the stated order, as assembled here from the cited components.
\end{proof}

\begin{theorem}[Conditional finite-ball equivalence]
\label{thm:ball}
Under Condition \ref{cond:NL}, the first law
$\delta S_B=\delta\langle H_B^{(\beta)}\rangle$ for all balls is
equivalent, at joint order $O(\beta\lambda)$, to the linearized EC
equations with torsionful boundary condition, where the bulk accounting
includes (i) the bulk axial field with boundary value $\beta$, (ii) the
metric response to the bilinear source $T[\,A(\beta),\delta\phi\,]$, and
(iii) the FLM bulk-entanglement cross term \cite{FLM2013}. Without
Condition \ref{cond:NL}, equality of the two sides fails by exactly
$\langle 0|\delta K_B^{\mathrm{nl}}|\delta\Psi\rangle+\mathrm{c.c.}$
\end{theorem}

\begin{proof}[Derivation]
Theorem 2.1 of \cite{McGwier2026} ($\delta S=\delta\langle K\rangle$) is
theory-independent and survives the deformation; what requires proof is
converting $\delta\langle K_B^{(\beta)}\rangle$ into bulk data. Splitting
$K_B^{(\beta)}=K_B+\delta K_B^{\mathrm{loc}}+\delta K_B^{\mathrm{nl}}$, the
first two terms are local integrals of $T_{00}$ and $\beta\cdot J_5$ with
known weights and map to the bulk by the standard dictionary, reproducing
the EC matching exactly as in the null case; the third term is the stated
obstruction. \pP{} as a reduction; Condition \ref{cond:NL} itself:
verified in the null-plane limit by Theorem \ref{thm:nullplane}, \pO{} for
finite balls.
\end{proof}

\begin{theorem}[Counterexample sector]
\label{thm:counterexample}
For matter for which the vector or mixed-symmetry parts of the spin current
do not vanish (non-minimal couplings, or non-Dirac matter with
non-conserved spin-current components), the corresponding components of
$b_{\lambda\mu\nu}$ source operators of dimension $d-1$ that are not
conserved currents. The deformation is then \emph{relevant}: it initiates
renormalization-group flow at first order, the deformed theory possesses no
conformal frame, the CHM construction of ball modular Hamiltonians is
inapplicable, and the naive torsionful first law---``first law for all
balls $\Leftrightarrow$ linearized EC with torsion boundary data''---is
false as stated in that sector.
\end{theorem}

\begin{proof}[Derivation]
A source coupled to a non-conserved operator of dimension $\Delta=d-1<d$
is a relevant coupling of dimension one; unlike the conserved-current case
there is no Ward identity protecting the coupling from renormalization,
and correlators acquire logarithmic flow at $O(b)$. The FGHMV equivalence
requires the CHM local modular Hamiltonian, which exists only in a CFT; a
flowing theory has none. The failure is structural, not computational.
\pP{} at the level of the dimensional and symmetry argument given; the
construction of an explicit non-minimal model exhibiting the failure
quantitatively is \pO{}.
\end{proof}

\section{Resolution of the Remainder}

Remark 6.3.3(i) of \cite{McGwier2026} is resolved as follows.

\begin{enumerate}[itemsep=2pt]
\item \textbf{Theorem sector.} For minimally coupled Dirac matter, torsion
sources are axial sources (Lemma \ref{lem:reduction}); pure-gauge sources
are trivial and only $d\beta$ acts locally (Lemma \ref{lem:gauge}); the
torsionful first law holds unconditionally on null planes
(Theorem \ref{thm:nullplane}) and holds for finite balls if and only if
the nonlocal modular remainder decouples (Theorem \ref{thm:ball},
Condition \ref{cond:NL}).
\item \textbf{Counterexample sector.} For generic torsion sources acting on
non-minimal matter, the deformation is relevant and the naive statement is
false (Theorem \ref{thm:counterexample}). The torsionful first law is
therefore \emph{not} universal; it is a property of the
conserved-current (axial/minimal) sector.
\item \textbf{Remaining open items.} \pO{}: (a) proof or refutation of
Condition \ref{cond:NL} for finite balls---the natural attack is the
S\'arosi--Ugajin kernel evaluated on the axial current's OPE block;
(b) joint order $O(\beta^2\lambda)$, where the Nieh--Yan anomaly channel
and the Hehl--Datta term both activate; (c) any statement for propagating
torsion.
\end{enumerate}

\section*{Certification}

Lemmas \ref{lem:reduction} and \ref{lem:gauge} are derived in full above.
Theorems \ref{thm:nullplane}--\ref{thm:counterexample} are assembled from
components each of which is proved in the primary literature cited at the
point of use, with the assembly performed here; the single unproved input
for the finite-ball case is isolated as Condition \ref{cond:NL} and never
used silently. The kernel normalization of perturbative modular theory is
deliberately not relied upon. Citations verified against the published
record during preparation: \cite{FLPW2016,SarosiUgajin2018,
CasiniTesteTorroba2017,DLP2017}. No claim is certified beyond its
derivation.

\section*{Acknowledgments}
The author acknowledges the use of Claude (Anthropic, claude.ai) for
assistance in drafting, literature verification, and \LaTeX{} preparation.
All technical claims and epistemic classifications were reviewed and
verified by the author, who bears sole responsibility for the content.

\begin{thebibliography}{99}\small

\bibitem{McGwier2026}
R.~W.~McGwier, ``Entanglement as the origin of spacetime curvature: a
verified synthesis of the first-law derivation of the Einstein equations,''
Version 3, Zenodo (2026), DOI:
\href{https://doi.org/10.5281/zenodo.21821890}{10.5281/zenodo.21821890}.

\bibitem{HehlDatta1971}
F.~W.~Hehl and B.~K.~Datta, ``Nonlinear spinor equation and asymmetric
connection in general relativity,''
J.\ Math.\ Phys.\ \textbf{12}, 1334 (1971).

\bibitem{HehlRMP1976}
F.~W.~Hehl, P.~von der Heyde, G.~D.~Kerlick and J.~M.~Nester, ``General
relativity with spin and torsion: foundations and prospects,''
Rev.\ Mod.\ Phys.\ \textbf{48}, 393 (1976).

\bibitem{NiehYan1982}
H.~T.~Nieh and M.~L.~Yan, ``An identity in Riemann--Cartan geometry,''
J.\ Math.\ Phys.\ \textbf{23}, 373 (1982).

\bibitem{ChandiaZanelli1997}
O.~Chand\'ia and J.~Zanelli, ``Topological invariants, instantons, and the
chiral anomaly on spaces with torsion,''
Phys.\ Rev.\ D \textbf{55}, 7580 (1997), arXiv:hep-th/9702025.

\bibitem{CHM2011}
H.~Casini, M.~Huerta and R.~C.~Myers, ``Towards a derivation of holographic
entanglement entropy,'' JHEP \textbf{05}, 036 (2011), arXiv:1102.0440.

\bibitem{FLM2013}
T.~Faulkner, A.~Lewkowycz and J.~Maldacena, ``Quantum corrections to
holographic entanglement entropy,'' JHEP \textbf{11}, 074 (2013),
arXiv:1307.2892.

\bibitem{Faulkner2014}
T.~Faulkner, M.~Guica, T.~Hartman, R.~C.~Myers and M.~Van Raamsdonk,
``Gravitation from entanglement in holographic CFTs,''
JHEP \textbf{03}, 051 (2014), arXiv:1312.7856.

\bibitem{FLPW2016}
T.~Faulkner, R.~G.~Leigh, O.~Parrikar and H.~Wang, ``Modular Hamiltonians
for deformed half-spaces and the averaged null energy condition,''
JHEP \textbf{09}, 038 (2016), arXiv:1605.08072.

\bibitem{CasiniTesteTorroba2017}
H.~Casini, E.~Test\'e and G.~Torroba, ``Modular Hamiltonians on the null
plane and the Markov property of the vacuum state,''
J.\ Phys.\ A \textbf{50}, 364001 (2017), arXiv:1703.10656.

\bibitem{DLP2017}
R.~Dey, S.~Liberati and D.~Pranzetti, ``Spacetime thermodynamics in the
presence of torsion,'' Phys.\ Rev.\ D \textbf{96}, 124032 (2017),
arXiv:1709.04031.

\bibitem{SarosiUgajin2018}
G.~S\'arosi and T.~Ugajin, ``Modular Hamiltonians of excited states, OPE
blocks and emergent bulk fields,'' JHEP \textbf{01}, 012 (2018),
arXiv:1705.01486.

\bibitem{GallegosGursoy2020}
A.~D.~Gallegos and U.~G\"ursoy, ``Holographic spin liquids and Lovelock
Chern--Simons gravity,'' JHEP \textbf{11}, 151 (2020), arXiv:2004.05148.

\bibitem{GGY2021}
A.~D.~Gallegos, U.~G\"ursoy and A.~Yarom, ``Hydrodynamics of spin
currents,'' SciPost Phys.\ \textbf{11}, 041 (2021), arXiv:2101.04759.

\bibitem{McGwier2026c}
R.~W.~McGwier, ``The New Standard Model for unified physics,''
Zenodo (2026), DOI:
\href{https://doi.org/10.5281/zenodo.21829536}{10.5281/zenodo.21829536}.

\bibitem{McGwier2026d}
R.~W.~McGwier, ``Condition NL for finite balls: a spectral reduction, the unconditional metric sector, and a universal kinematic residue,''
Zenodo (2026), DOI:
\href{https://doi.org/10.5281/zenodo.21836572}{10.5281/zenodo.21836572}.

\bibitem{McGwier2026e}
R.~W.~McGwier, ``The Hehl--Datta coefficient from relative-entropy positivity,''
Zenodo (2026), DOI:
\href{https://doi.org/10.5281/zenodo.21836757}{10.5281/zenodo.21836757}.

\bibitem{McGwier2026f}
R.~W.~McGwier, ``Emergent gravity and a new Lagrangian: a theory of everything we know,''
Zenodo (2026), DOI to be assigned.

\end{thebibliography}

\end{document}
